\begin{textsample}{POS Dim 6 – human – Score 3.00 – mg/mg\_ef\_linguagens\_ed\_fis\_2\_p12.txt}  \label{ex:f6_pos_020}
A Educação Física como componente curricular, tempo e espaço privilegiado de aprimoramento das dimensões motora/corporal, mental/cognitiva, afetivo/emocional, social/cultural e espiritual/existenciais, deverá ser trabalhada e desenvolvida no sentido de colaborar com a formação \textbf{integral} dos estudantes.
A partir deste documento curricular e da ampla discussão que deve ser realizada pela comunidade escolar, definindo suas singularidades, necessidades e possibilidades, é essencial que a construção dos seus documentos norteadores, entre eles o Projeto Político Pedagógico, seja feita de forma \textbf{coletiva}, valorizando práticas interdisciplinares e que contemple objetivos contemporâneos e estratégias formativas inovadoras.

% matched lemmas: coletivo, integral
\end{textsample}
